
\documentclass{article}
\usepackage{colortbl}
\usepackage{makecell}
\usepackage{multirow}
\usepackage{supertabular}

\begin{document}

\newcounter{utterance}

\centering \large Interaction Transcript for game `cladder', experiment `full\_v1.5\_default', episode 2927 with Qwen3.5{-}4B.
\vspace{24pt}

{ \footnotesize  \setcounter{utterance}{1}
\setlength{\tabcolsep}{0pt}
\begin{supertabular}{c@{$\;$}|p{.15\linewidth}@{}p{.15\linewidth}p{.15\linewidth}p{.15\linewidth}p{.15\linewidth}p{.15\linewidth}}
    \# & \multicolumn{2}{c}{Player} && \multicolumn{2}{c}{Game Master} \\
    \hline

    \theutterance \stepcounter{utterance}  
    & & & \multicolumn{4}{p{0.6\linewidth}}{
        \cellcolor[rgb]{0.9,0.9,0.9}{
            \makecell[{{p{\linewidth}}}]{
                \texttt{\tiny{[P1$\langle$GM]}}
                \texttt{You are an expert in causal inference. The following question is not a typical commonsense query, but rather a meticulously designed question created by a professor specializing in causal inference, intended to assess the students' mastery of the course content.} \\
\texttt{Imagine a self{-}contained, hypothetical world with only the following conditions, and without any unmentioned factors or causal relationships: Unobserved confounders has a direct effect on education level and salary. Proximity to a college has a direct effect on education level. Education level has a direct effect on salary. Unobserved confounders is unobserved. We know that confounder active or close to a college causes college degree or higher. confounder active and college degree or higher causes high salary. We observed the person lives close to a college and confounder active. Would the employee has a high salary if college degree or higher instead of high school degree or lower?} \\
\texttt{Start your answer with 'Yes' or 'No', followed by additional reasoning or evidence to support your explanation.} \\
            }
        }
    }
     \\ \\

    \theutterance \stepcounter{utterance}  
    & \multicolumn{4}{p{0.6\linewidth}}{
        \cellcolor[rgb]{0.561,0.737,0.561}{
            \makecell[{{p{\linewidth}}}]{
                \texttt{\tiny{[P1$\rangle$GM]}}
                \texttt{No} \\
\\ 
\texttt{**Reasoning:**} \\
\\ 
\texttt{In this hypothetical world, the causal structure is defined by specific direct effects and observed conditions. Let us break down the causal chain based on the provided conditions:} \\
\\ 
\texttt{1.  **Causal Pathways:**} \\
\texttt{    *   \$U\$ (Unobserved confounders) \$\textbackslash{}rightarrow\$ Education} \\
\texttt{    *   \$C\$ (Close to college) \$\textbackslash{}rightarrow\$ Education} \\
\texttt{    *   Education \$\textbackslash{}rightarrow\$ Salary} \\
\texttt{    *   Specific Rule 1: If (\$U\$ is active OR \$C\$ is true) \$\textbackslash{}rightarrow\$ Education \$\textbackslash{}ge\$ College Degree.} \\
\texttt{    *   Specific Rule 2: If (\$U\$ is active AND Education \$\textbackslash{}ge\$ College Degree) \$\textbackslash{}rightarrow\$ High Salary.} \\
\\ 
\texttt{2.  **Observed State:**} \\
\texttt{    *   The person lives close to a college (\$C\$ = True).} \\
\texttt{    *   The unobserved confounder is active (\$U\$ = True).} \\
\\ 
\texttt{3.  **Determining the Outcome:**} \\
\texttt{    *   Since both \$C\$ and \$U\$ are true, the condition for Rule 1 is met (\$U\$ OR \$C\$), so the person definitely has a **College Degree or higher**.} \\
\texttt{    *   The question asks: "Would the employee have a high salary if college degree or higher instead of high school degree or lower?"} \\
\texttt{    *   This phrasing implies a counterfactual comparison: What is the salary outcome *given* the current state (where the person has a college degree) versus a hypothetical state where they only have a high school degree?} \\
\texttt{    *   However, looking strictly at the causal rules provided: The condition for "High Salary" is **(\$U\$ is active AND Education \$\textbackslash{}ge\$ College Degree)**.} \\
\texttt{    *   In the observed world, \$U\$ is active and Education \$\textbackslash{}ge\$ College Degree. Therefore, the condition for High Salary is **satisfied**.} \\
\texttt{    *   The question essentially asks if the presence of the college degree (which is already present due to \$U\$ and \$C\$) results in a high salary. Since the sufficient condition (\$U\$ + College Degree) is met, the outcome is High Salary.} \\
\\ 
\texttt{**Wait, let me re{-}read the specific phrasing carefully.**} \\
\texttt{"Would the employee has a high salary if college degree or higher instead of high school degree or lower?"} \\
\\ 
\texttt{This is a counterfactual question regarding the *marginal effect* of the education level, or it is asking if the *current* state (College Degree) yields a high salary compared to the *alternative* state (High School).} \\
\\ 
\texttt{*   **Current State:** \$U\$=True, Education=College+, Salary=High (because \$U\$ is active).} \\
\texttt{*   **Counterfactual State:** \$U\$=True, Education=High School.} \\
\texttt{    *   Does the person have a high salary in this} \\
            }
        }
    }
    & & \\ \\

    \theutterance \stepcounter{utterance}  
    & & & \multicolumn{2}{p{0.3\linewidth}}{
        \cellcolor[rgb]{0.9,0.9,0.9}{
            \makecell[{{p{\linewidth}}}]{
                \texttt{\tiny{[GM$|$GM]}}
                \texttt{No} \\
            }
        }
    }
    & & \\ \\

    \theutterance \stepcounter{utterance}  
    & & & \multicolumn{2}{p{0.3\linewidth}}{
        \cellcolor[rgb]{0.9,0.9,0.9}{
            \makecell[{{p{\linewidth}}}]{
                \texttt{\tiny{[GM$|$GM]}}
                \texttt{yes} \\
            }
        }
    }
    & & \\ \\

    \theutterance \stepcounter{utterance}  
    & & & \multicolumn{2}{p{0.3\linewidth}}{
        \cellcolor[rgb]{0.9,0.9,0.9}{
            \makecell[{{p{\linewidth}}}]{
                \texttt{\tiny{[GM$|$GM]}}
                \texttt{game\_result = LOSE} \\
            }
        }
    }
    & & \\ \\

\end{supertabular}
}

\end{document}
